\section{Smoothing does not repair the canonical sharp interpolant (NS-62)}
\label{sec:ns62-sharp-interpolant}

\paragraph{Status and scope.}
The previous section changes the approximation norm and proves a smaller
upper budget. It is natural to ask whether simply smoothing the canonical
M\"obius interpolant now supplies convergence. It does not. The following
localized dual test extends the classical sharp-interpolant obstruction
in~\cite[Section 4]{BaezDuarte2000Arithmetic} through Mellin integration.
The conclusion concerns this explicit coefficient rule, including its
tail correction. It does not concern the optimized coefficients in
Proposition~\ref{prop:ns61-gain}.

Define
\[
 M(x)=\sum_{k\leq x}\mu(k),\quad
 s_N=\sum_{k\leq N}\frac{\mu(k)}k,\quad
 \eta_N=s_N-\frac{M(N)}N
       =\sum_{k<N}\frac{M(k)}{k(k+1)},
\]
and let
\begin{equation}
 B_N=\sum_{k=1}^N\mu(k)\rho_k-Ns_N\rho_N.
 \label{eq:ns62-natural-interpolant}
\end{equation}
The sign is chosen to match the classical notation: $B_N=-\chi$
on $x>1/N$ (up to endpoints). The sum is zero on $x>1$.
The smoothed target error is $J^r(\chi+B_N)$.

\begin{proposition}[A localized witness survives any number of integrations]
\label{prop:ns62-smoothed-obstruction}
For every integer $N\geq1$ and integer $r\geq0$,
\begin{equation}
 \norm{J^r(\chi+B_N)}_H
 \geq\frac{|N\eta_N-k_{N,r}|}{\sqrt N}
 \geq\frac{|N\eta_N-2|}{\sqrt N}
            -\frac{4\pi^2}{3^{r+2}N^{5/2}},
 \label{eq:ns62-laguerre-lower}
\end{equation}
where
\[
 k_{N,r}=\frac N{r!}\int_0^{1/N}
       \frac{\sin(2\pi x)}{\pi x}
                    \left(\log\frac1{Nx}\right)^r dx.
\]
In particular, for every sequence of nonnegative integers $r_N$,
$\norm{J^{r_N}(\chi+B_N)}_H$ does not tend to zero.
For growing $r_N$ the target also changes with $N$; this assertion
is only a failure of these particular approximants, not a new RH
criterion for varying norms.
The nonconvergence is for the full sequence of integers $N$. No
exclusion of a selected subsequence, including the dyadic subsequence,
is asserted without a corresponding lower bound for $\sqrt N|\eta_N|$
on that subsequence.
\end{proposition}
\begin{proof}
Use the dilation-invariant unitary $\mathcal V$ whose critical-line
Mellin multiplier, away from the discrete zero set, is
\[
 v(s)=\frac{s\zeta(1-s)}{(1-s)\zeta(s)}.
\]
Its modulus is one on $\Re s=1/2$; define it arbitrarily with modulus
one on the measure-zero exceptional set. Mellin identities and the
functional equation give
\[
 \mathcal V\rho_k(x)=\frac{\{kx\}}{kx},\qquad
 \mathcal V\chi(x)=k(x):=\frac{\sin(2\pi x)}{\pi x}.
\]
These are also the defining identities of the classical unitary in
\cite[Section 4]{BaezDuarte2000Arithmetic}. Both $\mathcal V$ and $J$
are Mellin multipliers, so they commute. For $0<x<1/N$,
\begin{equation}
 \mathcal V(\chi+B_N)(x)=k(x)+M(N)-Ns_N=k(x)-N\eta_N.
 \label{eq:ns62-unitary-small-cell}
\end{equation}

Let $L_r(z)=\sum_{j=0}^r(-1)^j\binom rj z^j/j!$ be the ordinary
Laguerre polynomial and put
\[
 \psi_{N,r}(x)=\sqrt N\,
       L_r\left(\log\frac1{Nx}\right)\mathbf1_{(0,1/N)}(x).
\]
Laguerre orthogonality \cite[Table 18.3.1]{DLMFLaguerre} gives
$\norm{\psi_{N,r}}=1$. The adjoint is
$(J^*f)(x)=x^{-1}\int_0^x f(y)\,dy$. Repeated integration, or
Rodrigues' formula and integration by parts, gives the complete
identity
\begin{equation}
 (J^*)^r\psi_{N,r}(x)=
 \frac{(-1)^r\sqrt N}{r!}
       \left(\log\frac1{Nx}\right)^r\mathbf1_{(0,1/N)}(x).
 \label{eq:ns62-adjoint-laguerre}
\end{equation}
There is no exterior term: at the $j$th integration, $j\leq r$,
the potential exterior coefficient is a Laguerre moment of a
polynomial of degree less than $r$, hence zero. On the interior,
the polynomial identity follows by integrating monomials against
$e^{-z}$; equivalently $(1-\partial_z)^{-r}L_r(z)=(-1)^rz^r/r!$.
This also covers $r=0$ directly.

Taking the inner product with $\psi_{N,r}$, moving $J^r$ to the
adjoint, and using~\eqref{eq:ns62-unitary-small-cell} yields
\[
 \left|\left\langle J^r\mathcal V(\chi+B_N),\psi_{N,r}\right\rangle\right|
       =\frac{|k_{N,r}-N\eta_N|}{\sqrt N}.
\]
The unitary norm and Cauchy--Schwarz prove the first inequality.
For real $x$, $|k(x)-2|\leq4\pi^2x^2/3$. Since
\[
 \int_0^{1/N}x^2\left(\log\frac1{Nx}\right)^r dx
        =\frac{r!}{3^{r+1}N^3},
\]
we have $|k_{N,r}-2|\leq4\pi^2/(3^{r+2}N^2)$, proving the second.

For completeness, $\eta_N\ne o(N^{-1/2})$ follows from the known
existence of a critical-line zeta zero. Indeed let
$\eta(x)=\int_1^xM(t)t^{-2}\,dt$. For $\Re s>0$,
\[
 \int_1^\infty\eta(x)x^{-s-1}\,dx
        =\frac1{s(s+1)\zeta(s+1)}.
\]
The right side has a genuine pole at $s=\rho-1$ for any
critical-line zero $\rho$. If $\eta(x)=o(x^{-1/2})$, the integral
would be analytic on $\Re s>-1/2$, and the usual Abelian boundary
estimate would make $(\sigma+1/2)$ times its value at
$\sigma+i\Im\rho$ tend to zero as $\sigma\downarrow-1/2$.
This contradicts a pole of any positive order at that point.
Also $\eta(N)=\eta_N$ and
$\eta(x)-\eta(\lfloor x\rfloor)=O(1/x)$, so the discrete
little-$o$ assertion would imply the continuous one. Thus
$\limsup_N\sqrt N|\eta_N|>0$, whereas the error term in
\eqref{eq:ns62-laguerre-lower} tends to zero uniformly in $r\geq0$.
This proves the stated nonconvergence for every choice of $r_N$.
No assumption about zeros off the critical line was made.
\end{proof}

\begin{proposition}[Positive smoothed differences still require signed coefficients]
\label{prop:ns62-positive-cone}
Every $\epsilon_n$, $n\geq2$, is nonnegative. Nevertheless $\tau$
is not in the $H$-closure of the cone of finite nonnegative linear
combinations of $\sigma_1,\epsilon_2,\epsilon_3,\ldots$.
This is a restriction on that coefficient cone, not on their signed
linear span or on the Weil form.
\end{proposition}
\begin{proof}
Nonnegativity was proved in Proposition~\ref{prop:ns61-average-budget}.
Its average identity also gives strict positivity on
$(1/n,1/(n-1))$.

Triangular telescoping gives
\begin{equation}
 -JB_N=\sum_{n=2}^{N}n s_{n-1}\epsilon_n.
 \label{eq:ns62-pointwise-series}
\end{equation}
It equals $\tau$ on $x>1/N$, because $-B_N=\chi$ there and
$Jf(x)$ uses only values at arguments at least $x$.
For $N=6$ the coefficient of $\epsilon_6$ is
$6s_5=6(1-1/2-1/3-1/5)=-1/5$.

Restrict now to $x>1/6$. Every $\epsilon_n$ with $n\geq7$
vanishes there. The restrictions of
$\sigma_1,\epsilon_2,\ldots,\epsilon_6$ are linearly independent:
first use $x>1$ for $\sigma_1$, then the successive intervals
$(1/n,1/(n-1))$ for $n=2,\ldots,6$.
Their finite-dimensional nonnegative cone is closed, but the unique
representation of the restricted target has the negative coefficient
$-1/5$. Thus the restricted target has positive distance from this
cone. Restriction is contractive, proving the full-space claim.
\end{proof}

\paragraph{Consequences for the next attempt.}
The positive shapes of the smoothed differences do not make the
coefficient problem positive. The canonical pointwise interpolation
coefficients are explicit, but their partial sums still fail in norm
after smoothing. The optimized residual criterion from NS-61 remains
open and is unaffected by this exclusion. A next attempt must obtain
an independent estimate for those adaptive residuals, or a different
coefficient rule with a complete norm bound; visual smoothness,
pointwise exactness away from zero, and a small raw block norm do not
supply it.
